% SHARD-ID: CA-44-QUBIT-SDP-EXCLUSION-HIERARCHY
% SHARD-TITLE: SDP Exclusion Hierarchy for Qubit Hamiltonians
% SHARD-SUMMARY: Packages finite Pauli moments, positivity, translation invariance, stationarity, and residual relations into a monotone SDP hierarchy.
% SHARD-SUMMARY: Infeasibility at any level excludes a fixed Hamiltonian density for the named generator route; feasibility proves nothing.
% SHARD-SUMMARY: Identifies which parameters must be fixed or scanned so the problem remains an SDP.
% SHARD-KEYWORDS: SDP hierarchy, exclusion, qubit Hamiltonian, moments, Poincare residuals, Virasoro residuals

\section{SDP Exclusion Hierarchy for Qubit Hamiltonians}
\label{sec:qubit-sdp-exclusion-hierarchy}

\subsection{Status, Sources, and Dependencies}

\textbf{Status: rigorous proposal.}  The hierarchy is a finite-state
consequence of the sourced state/GNS framework and the local residual
derivations.  No solver is implemented yet.

Dependencies: CA-38--CA-43, CONVENTIONS.md (n).

% Source: literature/md/2010.11121/2010.11121.md:178--184 -- projectively
%   consistent states and their projective limit.
% Source: literature/md/2010.11121/2010.11121.md:226--243 -- vacuum invariance
%   and warning that Poincare covariance is additional.

\subsection{Level Data}

Fix a Hamiltonian coefficient matrix \(h_{\alpha\beta}\).  Also fix, or scan
externally, the scalar subtraction \(e\), speed \(v^2\), and any coboundary
witnesses \(u,w\) used in CA-39.  This is important: if those witnesses are
unknown variables multiplied by unknown moments, the problem is no longer an
SDP without an additional lifting relaxation.

At level \(N\), choose finite sets
\[
  \mathcal W_N=\{\text{Pauli words supported in }[-N,N]\},
\]
\[
  \mathcal A_N=\{\text{probe words whose residual products fit the window}\},
\]
and residual densities
\[
  \mathcal R_N(h)=
  \{A-Du,\ B-u-v^2\bar h-Dw\}
\]
plus any selected finite Witt/Virasoro residuals from CA-42.

\subsection{SDP Constraints}

The variables are finitely many moments \(y_s\).  The constraints are
\[
  y_{\emptyset}=1,\qquad y_{\tau s}=y_s,
  \tag{44.1}
\]
\[
  M_{\mathcal W_N}(y)\succeq0,
  \tag{44.2}
\]
and optional stationarity constraints for local probes:
\[
  \omega(\delta_H(X))=0,\qquad
  \omega(\delta_P(X))=0,
  \qquad X\in\mathcal A_N.
  \tag{44.3}
\]
The generator-relation constraints are
\[
  \omega(X^*RY)=0,
  \qquad
  R\in\mathcal R_N(h),\quad X,Y\in\mathcal W_N,
  \tag{44.4}
\]
whenever the Pauli product \(X^*RY\) fits inside the moment support.  Each
entry in (44.1)--(44.4) is affine in \(y\) for fixed residual coefficients, so
this is a semidefinite feasibility problem.

\subsection{Exclusion Statement}

If a translation-invariant global state \(\omega\) exists whose GNS
representation carries the chosen local residual relations, then every finite
restriction of \(\omega\) satisfies every level \(N\).  Therefore:
\[
  \text{SDP level }N\text{ infeasible}
  \quad\Longrightarrow\quad
  h_{\alpha\beta}\text{ is excluded for this generator route.}
  \tag{44.5}
\]
The hierarchy is monotone: increasing \(N\) adds moment variables and relation
tests, and cannot turn a previously infeasible level into evidence of symmetry.

\subsection{Non-Implications}

Feasibility at all checked levels does not prove that a global state exists,
does not prove uniqueness or clustering, does not prove positive energy, and
does not prove a continuum Poincare or Virasoro representation.  It only means
that the tested finite restrictions have not yet found a contradiction.

When \(h_{\alpha\beta}\) is also a variable, the residual coefficients are
polynomial in \(h\).  Then the feasibility problem becomes a polynomial moment
problem over both Hamiltonian and state variables.  That is a separate
relaxation; the first implementation should instead take \(h\) as input and
return either an SDP infeasibility certificate or ``not excluded at this
level''.

